% =========================================================================
% Figura: Esquema simple del funcionamiento de un Autoencoder Variacional (VAE)
% Basado fielmente en el diagrama de referencia canónico:
% Input (x) -> Encoder -> [mu, sigma] -> z -> Decoder -> Reconstructed Input (x')
% Paleta de colores accesible Okabe-Ito definida en Configurations/01-Colors.sty
% =========================================================================

\begin{figure}[htbp]
\centering
\resizebox{0.92\linewidth}{!}{%
\begin{tikzpicture}[
    font=\sffamily,
    >=Stealth,
    x=1cm, y=1cm
]

% -------------------------------------------------------------------------
% 1. ENTRADA (INPUT) x
% -------------------------------------------------------------------------
\node[align=center, text=darkText] at (0.8, 2.15) {
    {\fontsize{11}{13}\selectfont\textbf{Entrada}}\\[1pt]
    {\footnotesize\color{grayText} (\textit{Input})}
};

\draw[draw=grayText!70, fill=cardBorder!25, rounded corners=5pt, line width=1.1pt]
    (0.0, -1.6) rectangle (1.6, 1.6);
\node[font=\fontsize{16}{18}\selectfont\bfseries, text=darkText] at (0.8, 0) {$\mathbf{x}$};

% Flecha Entrada -> Codificador
\draw[line width=1.1pt, draw=grayText!90!black, -{Stealth[length=5pt, width=3.5pt]}]
    (1.6, 0) -- (2.4, 0);

% -------------------------------------------------------------------------
% 2. CODIFICADOR (ENCODER)
% -------------------------------------------------------------------------
\filldraw[
    draw=figBlue,
    fill=figBlueTint,
    rounded corners=5pt,
    line width=1.2pt
] (2.4, -1.7) -- (5.2, -1.0) -- (5.2, 1.0) -- (2.4, 1.7) -- cycle;

\node[align=center, text=darkText] at (3.8, 0) {
    {\fontsize{11}{13}\selectfont\textbf{\color{figBlue!90!black} Codificador}}\\[3pt]
    {\footnotesize\color{grayText} (\textit{Encoder})}
};

% Conexión bifurcada de Codificador a mu y sigma
\draw[line width=1.1pt, draw=grayText!90!black] (5.2, 0) -- (5.8, 0);
\draw[line width=1.1pt, draw=grayText!90!black] (5.8, -0.9) -- (5.8, 0.9);
\draw[line width=1.1pt, draw=grayText!90!black, -{Stealth[length=4.5pt, width=3pt]}] (5.8, 0.9) -- (6.4, 0.9);
\draw[line width=1.1pt, draw=grayText!90!black, -{Stealth[length=4.5pt, width=3pt]}] (5.8, -0.9) -- (6.4, -0.9);

% -------------------------------------------------------------------------
% 3. PARÁMETROS LATENTES: \mu Y \sigma
% -------------------------------------------------------------------------
% Caja \mu
\draw[draw=figOrange, fill=figOrangeTint, rounded corners=4pt, line width=1.1pt]
    (6.4, 0.4) rectangle (8.2, 1.4);
\node[font=\fontsize{15}{17}\selectfont\bfseries, text=darkText] at (7.3, 0.9) {$\boldsymbol{\mu}$};

% Caja \sigma
\draw[draw=figOrange, fill=figOrangeTint, rounded corners=4pt, line width=1.1pt]
    (6.4, -1.4) rectangle (8.2, -0.4);
\node[font=\fontsize{15}{17}\selectfont\bfseries, text=darkText] at (7.3, -0.9) {$\boldsymbol{\sigma}$};

% Conexión convergente de mu y sigma a z
\draw[line width=1.1pt, draw=grayText!90!black] (8.2, 0.9) -- (8.8, 0.9);
\draw[line width=1.1pt, draw=grayText!90!black] (8.2, -0.9) -- (8.8, -0.9);
\draw[line width=1.1pt, draw=grayText!90!black] (8.8, -0.9) -- (8.8, 0.9);
\draw[line width=1.1pt, draw=grayText!90!black, -{Stealth[length=4.5pt, width=3pt]}] (8.8, 0) -- (9.4, 0);

% -------------------------------------------------------------------------
% 4. VECTOR LATENTE z
% -------------------------------------------------------------------------
\draw[draw=figOrange!90!black, fill=figOrangeTint, rounded corners=4pt, line width=1.2pt]
    (9.4, -0.6) rectangle (11.0, 0.6);
\node[font=\fontsize{15}{17}\selectfont\bfseries, text=darkText] at (10.2, 0) {$\mathbf{z}$};

% Ecuación de reparametrización debajo de z
\node[align=center, text=darkText] at (10.2, -1.65) {
    {\fontsize{11}{13}\selectfont $\mathbf{z} = \boldsymbol{\mu} + \boldsymbol{\sigma}\boldsymbol{\epsilon}$}\\[3pt]
    {\footnotesize\color{grayText} $\boldsymbol{\epsilon} \sim \mathcal{N}(\mathbf{0}, \mathbf{I})$}
};

% Flecha z -> Decodificador
\draw[line width=1.1pt, draw=grayText!90!black, -{Stealth[length=5pt, width=3.5pt]}]
    (11.0, 0) -- (11.8, 0);

% -------------------------------------------------------------------------
% 5. DECODIFICADOR (DECODER)
% -------------------------------------------------------------------------
\filldraw[
    draw=figGreen,
    fill=figGreenTint,
    rounded corners=5pt,
    line width=1.2pt
] (11.8, -1.0) -- (14.6, -1.7) -- (14.6, 1.7) -- (11.8, 1.0) -- cycle;

\node[align=center, text=darkText] at (13.2, 0) {
    {\fontsize{11}{13}\selectfont\textbf{\color{figGreen!80!black} Decodificador}}\\[3pt]
    {\footnotesize\color{grayText} (\textit{Decoder})}
};

% Flecha Decodificador -> Salida
\draw[line width=1.1pt, draw=grayText!90!black, -{Stealth[length=5pt, width=3.5pt]}]
    (14.6, 0) -- (15.4, 0);

% -------------------------------------------------------------------------
% 6. ENTRADA RECONSTRUIDA (RECONSTRUCTED INPUT) x'
% -------------------------------------------------------------------------
\node[align=center, text=darkText] at (16.2, 2.15) {
    {\fontsize{11}{13}\selectfont\textbf{Reconstrucción}}\\[1pt]
    {\footnotesize\color{grayText} (\textit{Output} $\mathbf{x}'$)}
};

\draw[draw=grayText!70, fill=cardBorder!25, rounded corners=5pt, line width=1.1pt]
    (15.4, -1.6) rectangle (17.0, 1.6);
\node[font=\fontsize{16}{18}\selectfont\bfseries, text=darkText] at (16.2, 0) {$\mathbf{x}'$};

\end{tikzpicture}%
}
\caption[Diagrama esquemático del funcionamiento de un Autoencoder Variacional (VAE)]{Diagrama esquemático del funcionamiento de un Autoencoder Variacional (VAE). Adaptado de~\cite{kim2024exploring}.}
\label{fig:esquema_vae}
\end{figure}
